\cotefile{14}
\scanpage{31}
2. \uline{Familles de classes de cycles \deleted{\illegible{1} \unsure{à équiv.|à équivalence}} \inserted{mod équiv. d'Alb.}}

\uline{Théorème} 2.1. Soient $X^n$ et $Y^m$ comme d'habitude,
et $Z\in\deleted{\unsure{\operatorname{Cyc}}}A^{i}_{\inserted{\mathrm{alb}}}(X\times Y)$, d'où une \deleted{\unsure{hom composé|hom composée|loi compo}} \inserted{application}
$\unsure{\mathfrak{a}|\mathfrak{z}|\mathfrak{t}}\leadsto Z(\mathfrak{a})$ de $\deleted{\unsure{A_0(Y)}}\inserted{A_{0\,\mathrm{alb}}(Y)}$ dans $A^i_{\mathrm{alb}}(X)$. \deleted{Il existe une
application et une seule de $(A_{0\,\mathrm{alb}}(Y)\simeq\mathrm{Alb}^{*}_{\deleted{Y}}(k))$ dans $A^{\unsure{+|i|*}}_{\mathrm{alb}}(X)$
telle que pour tout $\mathfrak{a}\in\mathrm{Alb}_0(Y)$,} or $A_{0\,\mathrm{alb}}(Y)\overset{S_Y}{\simeq}\deleted{\unsure{B}}A_Y^*(k)$,
d'où un hom de groupes
\[
u_Z^*:\underset{\inserted{\substack{\shortparallel\\\mathrm{Alb}}}}{A_Y}^{*}(k)\longrightarrow A^i_{\mathrm{alb}}(X)
\]
\deleted{dont la restriction à $A_Y^{(n)}(k)$ est} défini par la
condition
\[
\deleted{\unsure{S_Y}}Z(\mathfrak{a})=u_Z(S(\mathfrak{a}))
\]
pour tt $\mathfrak{a}$ (resp. pour tt $\mathfrak{a}$ de degré 0,
si on regarde la restriction en degré 0).

\uline{De plus}, si on regarde l'hom \uline{induit}
sur $A_Y(k)=A_Y^0(k)$
\[
u_Z:A_Y(k)\longrightarrow A^i_{\mathrm{alb}}(X),
\]
on a ce qui suit : \deleted{il existe une hom de
variétés de VA
\[
p:A'\longrightarrow A_Y
\]
et un cycle
\[
Z'\in A^i_{\mathrm{alb}}(\unsure{/}
\]
}
(i) Il ne dépend que de la classe de $Z$ dans
$A^{i\,\tau}_{\deleted{\unsure{\mathrm{alg}}}}(X\times Y)$ \margin{[A réserver pour plus tard}

(ii) Il existe une \deleted{\unsure{homo|base}}\inserted{isogénie}
% archive sheet number: 30
\scanpage{32}
de VA
\[
\varphi:A'\longrightarrow A_Y
\qquad\text{\raisebox{1ex}{$\checkmark$}}
\]
et un $Z'\in A^i_{\mathrm{alb}}(\deleted{\unsure{A'|X}}X\times A')$, \deleted{dont} tel que
$Z'((0_{A'}))=0$ \inserted{diviseur de degré 1 concentré en 0} (ou ce qui revient au $=$, tel que
l'application $v_{Z'}:\deleted{\unsure{Z}}a'\leadsto Z'(a')$ de $A'(k)$ dans $\deleted{\unsure{A'|A^i_Z}}A^i_{\mathrm{alb}}(X)$
soit un hom de groupes) tel que
\[
u_Z\circ\varphi(k)=v_{Z'}
\]

\uline{Dém} Supposons qu'il existe un cycle
$\deleted{\unsure{Z'|Z}}U$ sur $Y\times A_Y$, de codim $m$, tel
que pour tt $a\in A_Y$, on ait $S_Y(U(a))=a$.
Alors la ``famille algébrique de diviseurs''
composée $a\leadsto Z(U(a))=Z'(a)$
convient. \unsure{Il faut constater|On doit constater|Il faut voir} que pour tt $\unsure{\mathfrak{a}|\mathfrak{z}|\mathfrak{t}}\in\deleted{\unsure{A}}A_{0\,\mathrm{alb}}(\unsure{Y|X})$
\inserted{de degré 0,} on aura \deleted{$\unsure{S_Y Z'(a)|Z'(a)}$.}
\[
Z(\mathfrak{a})=Z'(S(\mathfrak{a}))
\]
Car pour le voir, on \deleted{peut \unsure{remplacer} \unsure{à priori|d'abord} $\mathfrak{a}$ par un cycle \unsure{Alb-équivalent}} \inserted{note que $\mathfrak{a}$ est de la forme} $U(a)$, \deleted{\unsure{donc|dont} il \unsure{faut}} vérifier
\[
\begin{aligned}
Z(U(a))&=Z'(S(U(a))\\
&\phantom{=}\shortparallel\\
&\phantom{=}Z'(a)
\end{aligned}
\qquad \text{O.K.}
\]

Mais \deleted{à priori,} on ne sait pas si on peut
% archive sheet number: 31
\scanpage{33}
trouver $U$ \inserted{comme dessus}, mais on peut trouver $U$ de
façon que l'on ait
\[
S_Y(U(a))=na
\]
d'où, posant $Z'(a)=Z(U(a))$, la relation
\[
Z'(S_Y(\unsure{\mathfrak{a}|\mathfrak{z}|\mathfrak{t}}))=nZ(\mathfrak{a})\qquad \mathfrak{a}\ \text{de degré 0 sur }Y.
\]
\deleted{Car on peut trouver $nS(\mathfrak{a})=U(a)$ si $a=S_Y(\mathfrak{a})$,
$S(\mathfrak{a})$ de la forme $U(\unsure{a|a'})$, car si $a'=S_Y(\mathfrak{a})$,
on aura $S_Y(U(a'))=na'=S_Y(n\mathfrak{a})$ donc
$U(a')=n\mathfrak{a}\unsure{)}$, et comme on aura $a'=na$ pour
$a$ convenable, il s'ensuivra $U(a)$}

Pour le voir, soit $a=S_Y(\mathfrak{a})$, on aura
$\deleted{\unsure{Z'(S}}S_Y(U(a))=na=S_Y(n\mathfrak{a})$ d'où \deleted{$U(a)=n\mathfrak{a}$,}
\[
U(a)=n\mathfrak{a}
\]
d'autre part on aura
\[
\begin{aligned}
\deleted{\unsure{S}}Z'(S_Y(U(a)))&=nZ(U(a))\\
\shortparallel\quad &\phantom{=}\shortparallel\\
Z'(na)\quad &\phantom{=}nZ'(a)\quad .
\end{aligned}
\]
\deleted{d'où, comme $U(a)=n\mathfrak{a}$}
Or tout \unsure{élément|élém} de $A_{0\,\mathrm{alb}}(Y)$ de degré 0 est
\deleted{$nZ'(S_Y(\mathfrak{a}))=n^2Z(\mathfrak{a})$}
de la forme $n\unsure{\mathscr L|\mathfrak b|L|\mathcal L}$ (car $A_Y(k)$ est divisible !)
donc appliquant ce qui précède à $\mathscr L$ au lieu de $\mathfrak{a}$,
on trouve
\[
Z'(S_Y(\mathfrak{a}))=nZ(\mathfrak{a}).
\]

On prend comme isogénie la multiplication par $n$.
% archive sheet number: 32
\scanpage{34}
\unsure{Cela|Ceci|On} prouve (ii). Pour (i), on note qu'on sait
déjà que si $Z$ est alg. équiv. à 0, alors
l'hom correspondant $u:A_Y(k)\longrightarrow A^i_{\mathrm{alb}}(X)$ est nul
(th. 1.6.), donc si $Z$ est $\tau$-équivalent à 0, alors
l'hom $u$ est tel que $mu=0$ pour quelque
$m>0$, donc $u=0$ car $A_Y(k)$ est $m$-divisible.

\uline{Corollaire} \unsure{1.7.|2.2.|2.7.} Pour qu'un cycle algébrique
de codim $i$ sur $X$ soit alg. équiv. à 0, il
f. et s. qu'il soit Albanese-équivalent à
un diviseur de la forme $Z(a_0)$, où
\deleted{$A$ est une} $Z$ est un cycle de \unsure{diodim|dicodim|dicdim|codim} $i$
sur un \deleted{VA $A$,} produit $\underline{A}\times X$, $A$ une
variété abélienne, tel que $\deleted{\unsure{A(k)|\mathrm{Alb}}}a\leadsto Z(a):$
$A\longrightarrow\deleted{\unsure{Z^i(}}A^i_{\mathrm{alb}}(X)$ soit un hom de
groupes.

\uline{Corollaire} 2.3. Pour tt entier $n>0$, le ss-groupe
$P^i(X)$ \unsure{des|de} \deleted{classes d'équivalence, à équiv. d'Alb.} \inserted{$A^i_{\mathrm{alb}}(X)$, formé des classes de cycles}
alg. équiv. à 0, est $n$ \uline{divisible}.

\uline{Corollaire} 2.4. Soient $Z,Z'\in A^*_{\mathrm{alb}}(X)$, dont l'un est
alg. équiv. à 0 et l'autre $\tau$-équiv. à 0.
Alors $ZZ'=0$.
% archive sheet number: 33
\scanpage{35}
\uline{Définition} 2.5. Soient $X,Y,\ldots$ et considérons une
application
\[
\varphi:\underline{Y(k)}\longrightarrow \tikz[baseline=(terme.base)]{\node[inner sep=1pt] (terme) {$A^i_{\mathrm{alb}}(X)$};\draw[gray,thin,rotate around={-5:(terme.center)}] (terme.center) ellipse (2.7em and .85em);}
\]
On dit que cette application est \uline{algébrique} si

$1^\circ)$ L'hom $\deleted{Z}\operatorname{Cyc}_0(Y)\longrightarrow A^i_{\mathrm{alb}}(X)$
correspondant se factorise par
\margin{\diagram{\begin{tikzcd}[ampersand replacement=\&,column sep=tiny,font=\scriptsize]
Y\arrow[r]\& A_Y^1\arrow[r,"N"]\& A_Y^N
\end{tikzcd}}}
$\deleted{\unsure{A_Y}}A_{0\,\mathrm{alb}}(Y)\simeq A_Y^*(k)$, i.e. \deleted{\unsure{que|s'annule}} est
nulle sur les $\unsure{\mathfrak{a}|\mathfrak{z}|\mathfrak{t}}$ de degré 0 tels
que $S_Y(\mathfrak{a})=0$

$2^\circ)$ L'hom induit
\[
\varphi:\tikz[baseline=(terme.base)]{\node[inner sep=1pt] (terme) {$A_Y^{\unsure{0|*}}(k)$};\draw[gray,thin,rotate around={-5:(terme.center)}] (terme.center) ellipse (2.7em and .85em);}\longrightarrow A^i_{\mathrm{alb}}(X)
\]
est tel que pour $N>0$ convenable, $\underline{N\varphi}(x)=\varphi(Nx)$
puisse être induit par un cycle algébrique
\[
Z\in A^i_{\mathrm{alb}}(X\times A_Y)
\]

\uline{Prop.} 2.6. Si $\varphi$ est définie par un $Z\in A^i_{\mathrm{alb}}(X\times Y)$,
elle est \uline{algébrique}.

C'est 2.1. Va être généralisé en 2.7. (ii)

\uline{Prop.} 2.7. \inserted{(0) Forment un \uline{ss-groupe}. Si $N\varphi$ est alg., $\varphi$ l'est.}
(i) Fonctorialité en $Y$. \deleted{Plus précisément}

(ii) Fonctorialité en $X$, plus précis, si on a
un hom $\deleted{P}A^i_{\mathrm{alb}}(X)\longrightarrow\deleted{P}A^{\unsure{j|i'}}_{\mathrm{alb}}(X')$ défini par
une classe de correspondance algébrique \uline{à coeff.
\deleted{\unsure{rationnels}}} entiers.
% archive sheet number: 34
\scanpage{36}
N.B. Fixons un pt $y_0$ de $Y$. Alors le groupe
des familles algébriques $Y(k)\unsure{-|\longrightarrow}A^i_{\mathrm{alb}}(X)$ est
isom à $A^i_{\mathrm{alb}}(X)\times\operatorname{Hom\,alg}(A_Y,A^i_{\mathrm{alb}}(X))$, où le
deuxième facteur désigne les familles algébriques
qui sont des hom de groupes, i.e. nulles à
l'origine. On peut donc \deleted{se} limiter à étudier
les hom ``alg.'' $A_{\deleted{\unsure{Y}}}(k)\longrightarrow A^i_{\mathrm{alb}}(X)$, \inserted{où $A$ est une VA,} et interprétant
$A$ comme $A_{0\,\mathrm{alb}}(A)=A^\alpha_{\mathrm{alb}}(A)^{(\text{alg.équiv.}\,0)}$ $(\alpha=\dim A)$
on peut considérer \deleted{\unsure{le}} ces objets comme des
cas particuliers des \uline{hom algébriques}
\margin{$A(k)\simeq P^{\unsure{\alpha|n}}(A)\longrightarrow P^i(X)$}
\[
\varphi:P^j(X')\longrightarrow P^i(X).
\]

Un hom de gpes $\varphi$ est dit ``algébrique''
s'il existe un entier $N>0$ tel que
$N\varphi$ soit définissable par un cycle algébrique sur $X\times\unsure{Y'|X'}$
(donc $\varphi$ serait définissable par un cycle alg.
à coeff. \uline{rationnels} --- mais attention, un cycle
à coeff. rationnels \uline{quelconque} ne définit pas en
général un hom $P^j(X')\longrightarrow P^i(X)$, $=$ si les
degrés sont les bons ....). Ces hom \uline{se composent}
% archive sheet number: 35
\scanpage{37}
\margin{$?\quad\big|$}Une question plus subtile est s'ils peuvent
s'inverser, quand c'est des isomorphismes ....

Dualité entre $P^i(X)$ et $P^{\deleted{n+1-i}\inserted{j}}(X)$ \inserted{pour $i+j=n+1$} : Si $A,B$
\deleted{\unsure{sont une|sont deux}} \inserted{des} VA, \deleted{dont la dualité est} accouplées
par une dualité parfaite $(A\simeq B^{\unsure{*|\vee}}, B\simeq A^*)$,
alors il y a correspondance biunivoque
entre $\operatorname{Hom\,alg}(\deleted{\unsure{B}}A,P^i(X))$ et $\operatorname{Hom}^{\mathrm{alg}}(P^{\unsure{1|j}}(X),B)$ :

En effet, si on a un hom
$A\longrightarrow P^i(X)$, défini par un $Z\in A^i_{\mathrm{alb}}(\deleted{A}X\times\deleted{X}A)$,
alors $ {}^t Z$ définit un hom $P^j(X)\longrightarrow P^{\deleted{\unsure{1}}i+j-n}(A)=\underset{\substack{\shortparallel\\B}}{P^1(A)}$
et inversement si on a $P^j(Y)\longrightarrow B$. \deleted{// Notation}
\deleted{$\unsure{{}^t\varphi|\mathrm{tq}}$ \unsure{s'impose}. Plus généralement}
Cependant, \margin{$\big|\big|\big|$}il faut vérifier que $P^j(X)\longrightarrow B$ ne dépend que
de \uline{l'application} $A\longrightarrow P^i(X)$, pas du cycle $Z$.

Question de transposition se pose, plus généralement, pour
$P^i(X)\unsure{-|\longrightarrow}P^{\unsure{i'|a'|j'}}(X')$ en $P^{j'}(X')\unsure{-|\longrightarrow}P^j(X)$.
Cf plus bas ....
% archive sheet number: 36
\scanpage{38}
3. Homomorphismes $\ell^{\unsure{b|*}}$-adiques associés à une
\deleted{\unsure{hom}} famille abélienne de cycles.

Soit $A$ une V.A., $Z\in A^i_{\mathrm{alb}}(X\times A)$, d'où
\[
\deleted{\unsure{\varphi=}}u_Z:A(k)\longrightarrow A^i_{\mathrm{alb}}(X)
\]
mais aussi
\[
a({}^t Z):H^{2(n-i)+1}(X)(n-i)\longrightarrow H^1(A)
\]
(coefficients $\ell$-adiques $\deleted{\unsure{\mathbb Q}}\mathbb Z_\ell$). Je dis que
mod torsion (i.e. à \unsure{coefficient|coeffient} $\mathbb Q_\ell$ !) cet
homomorphisme ne dépend que de $u_Z$, pas de $Z$,
\deleted{donc} et de façon précise que

\margin{? A-t-on une réciproque}
% The two pencil ovals around the heading and the marginal note are joined.
\uline{Théorème} 3.1. \margin{dém. pas faite}
\[
u_Z=0\ \overset{?}{\underset{\unsure{-|\sim}}{\unsure{\vert\!\Longrightarrow|\Longleftrightarrow|\Longrightarrow}}}\ a({}^t Z):H^{\unsure{2(n-i)+1|(n-i)+1}}\unsure{-|\to}H^1\ \text{est nul}
\]

\uline{Démonstration} Choisissons une ``courbe ample''
$C$ sur $A$, \deleted{\unsure{d'où alors|disons}} d'où $C\longrightarrow A$ induisant
$\deleted{C\times X\unsure{\longrightarrow|-} A_C\longrightarrow A}\inserted{X\times C\longrightarrow X\times A}$, d'où $Z'\in A^i_{\mathrm{alb}}(C\times X)$ et
par conséquent
\[
u_{Z'}:C(k)\longrightarrow A^i_{\mathrm{alb}}(X)
\]
\[
a({}^t Z'):H^{2(n-i)+1\deleted{\unsure{3}}}\deleted{(X)}(n-i)\longrightarrow H^1(C)
\]
et comme $H^1(A)\longrightarrow H^1(C)$ est injectif, $C(k)\longrightarrow A(k)$
surjectif, --- .
% archive sheet number: 37
\scanpage{39}
\[
\begin{aligned}
u_Z=0&\ \unsure{\Longleftrightarrow|\Longleftarrow}\ u_{Z'}=0\\
\deleted{\unsure{u}}a({}^t Z)=0&\ \Longleftrightarrow\ a({}^t Z')=0
\end{aligned}
\]
On est donc ramené à la question analogue,
avec $A$ remplacé par $C$. Or \deleted{tout
cycle} la décomposition
\[
H^N(X\times C)\simeq H^N(X)\otimes H^0(C)+H^{N-1}(X)\otimes H^1(C)+H^{N-\unsure{2|1}}(X)\otimes H^{\unsure{2|1}}(C)
\]
est algébrique, comme \deleted{\illegible{1}} on vérifie aisément.
Appliquant ceci au cas $N\unsure{=|\,}2i$, et notant
qu'un cycle algébrique appartenant à l'un des
deux facteurs extrêmes donne un $u_{Z'}$ \uline{et}
un $a({}^t Z)$ nul, on est ramené au
cas d'un cycle appartenant au facteur
médian, i.e. dont la \uline{projection} sur $X$, et
la ``restriction'' à $X$, sont \unsure{trivles|triviales}.
% archive sheet number: 38 (last digit overwritten; possibly 39)
\scanpage{40}
Astuce de Tate sur les cycles décomposables
% archive sheet number: 39
\scanpage{41}
Soit $A$ une VA sur $k$ alg. clos,
d'où module $\mathbb T_\ell(A)\deleted{\unsure{\otimes}}$ sur la pro-alg. de Lie $\deleted{\unsure{\mathfrak g_K}}\unsure{\mathfrak g|\mathscr G}(k)=\mathfrak g$.
Soit $K$ une clôture alg. de $\mathbb Q_\ell$, posant
\[
\mathbb T_{K}(A)=\mathbb T_\ell(A)
\]
de sorte que $\mathbb T_K(A)$ est un module sur
$\mathfrak g\otimes_{\mathbb Q_\ell}K=\mathfrak g_K$. Supposons

a) La \deleted{\unsure{hyp}} conj de Tate vérifiée pour $H^2(A,\mathbb Q_\ell(1))$
\[
H^2(A\deleted{)},\mathbb Q_\ell(1))^{\mathrm{alg}}=H^2(A,\mathbb Q_\ell(1))^{\mathfrak g}
\]

b) $\mathfrak g_K$ opérant sur $\deleted{\unsure{\mathfrak g}}\deleted{\text{est}}\mathbb T_K(A)$ est ``pseudoalgébrique''
[i.e. les valeurs propres des frobenius \deleted{sont} \unsure{de|des}
norme $q$ sont $\varepsilon q^{-1/2}$, $\varepsilon$ une racine de l'unité \ldots]

% slanted
\margin{\unsure{Par|Pour} Tits, le premier cas est un cas particulier du second}
\uline{ou} $\mathbb T_K(A)$ se décompose en deux morceaux
$V_1\times V_2$, tels que chacun soit propre pour ``presque
tous'' les Frobénius, avec des valeurs propres
différentes pour presque tous les Frobénius.

[Ces conditions sont remplies si $A$ est une \deleted{produit}\inserted{puissance $C^n$}
\deleted{de}\inserted{d'une} courbe elliptique \inserted{$C$ définissable} sur un corps fini $\mathbb F_q$,
comme on voit en notant que a) \deleted{\inserted{\unsure{résulte}} et b)
sont stables par passage à une puissance
cartésienne, et pour une $C$, a) est}
a été vérifié par Mumford (on est ramené
% archive sheet number: 40
\scanpage{42}
au cas du produit de \uline{deux} courbes
elliptiques), et b) est trivial à partir du
cas de $C$, car b) est stable par \deleted{produit}
puissance cartésienne]. Alors on a pour tt $i$
\[
\unsure{D|\mathcal D}^{2i}(A,\mathbb Q_\ell(i))=H^{2i}(A,\mathbb Q_\ell(i))^{\mathrm{alg}}=H^{2i}(A,\mathbb Q_\ell(i))^{\unsure{\mathfrak g|\mathscr G|U}}
\]
\margin{décomposables $\mathbb Q_\ell$-algébriques}

En effet, \deleted{l'hypothèse} à priori on a les
inclusions $\subset$, $\subset$, et l'hypothèse implique
l'égalité des extrêmes ....
% archive sheet number: 41
\scanpage{44}
27.8.65\\

Mon Cher Serre,\\

Je continue à retourner dans tous les sens les questions de cycles
algébriques; techniquement parlant je vois la question plus clairement, mais
cependant il manque toujours le pas décisif. Comme je te connais allergique
à la cohomologie, j'ai envie de t'énoncer les deux conjectures-clef, qu'il
s'agit de démontrer, de façon purement géométrique i.e. sans parler de
cohomologie. $X$ désigne une variété projective lisse connexe sur le corps
algébriquement clos $k$ , de dimension $n$, $Y$ une section hyperplane lisse.
Par $C^i(X)$ je désigne le groupe des classes de cycles de codimension $i$
mod équivalence algébrique, tensorisé par $\underline{\mathrm Q}$ , par $\handwritten{\xi}\ \handwritten{\in}\ C^1(X)$ la classe
de $Y$

\uline{Conjecture} A\qquad Pour tout entier $i$ tel que $2i\handwritten{\leqslant}n$, le produit par $\handwritten{\xi}^{n-2i}$
induit un isomorphisme $C^i(X)\handwritten{\xrightarrow{\sim}}C^{n-i}(X)$.

Modulo vérification que je n'ai pas écrite en détail, il suffit
même dans cette conjecture de vérifier la surjectivité, et de le faire pour
$n=2i+1$. Le premier cas qui échappe pour le moment est $i=1,n=3$.
Cette conjecture implique l'analogue, pour les $C^i$, des théorèmes connus
de Lefschetz sur la comparaison de la cohomologie de $X$ et de $Y$ par les
homomorphismes d'image directe et d'image inverse (pour la formulation
correcte, se rappeler seulement que $C^i$ devient $H^{2i}$), en même temps
qu'il implique les théorèmes de Lefschetz cohomologiques eux-mêmes, dont
la forme formellement la plus forte consiste à affirmer que $\handwritten{\xi}^{n-i}$ induit
un isomorphisme de $H^i$ et $H^{2n-i}$ pour la cohomologie à coefficients $\underline{\mathrm Q}_{\rlap{$\scriptstyle\handwritten{/}$}\unsure{\mathrm{l}|1|\ell}}$ ;
je te signale que ce théorème purement cohomologique n'est pas même démontré
pour le moment ! N'est démontré que le théorème de Lefschetz
\textquotedbl{}faible\textquotedbl{} affirmant que la dimension cohomologique d'une variété affine
de dimension $n$ est majorée par $n$, ce qui pour la variété affine
$U=X-Y$ prend la forme : $H^i(X)\handwritten{\longrightarrow}H^i(Y)$ est bijectif pour $i\handwritten{\leqslant}n-2$,
injectif pour $i=n-1$, ou encore : $H^i(Y)\handwritten{\longrightarrow}H^{i+2}(X)$ est bijectif pour

\hfill .../...
% archive sheet number: 43
